\chapter{Дефекты кристаллов}\label{ch:defects}

\section{Точечные дефекты}

Даже идеальный кристалл при $T>0$ содержит \textbf{вакансии} --- пустые узлы. Их равновесная доля
\begin{equation}
c_v = e^{S_f/k_B}\,e^{-H_f/k_BT},
\end{equation}
$H_f\sim1$--2~эВ, у точки плавления $c_v\sim10^{-4}$. Образование вакансии стоит энергии, но приносит
энтропию --- та же борьба $H$ и $TS$, что и в главе~\ref{ch:gibbs}. Атомы \textbf{внедрения} (водород, кислород, углерод)
сидят в пустотах между узлами. Именно поэтому водород не нуждается в вакансиях, чтобы двигаться.

\section{Дислокации}

Пластическая деформация металла идёт не сдвигом целых атомных плоскостей разом (на это нужны напряжения порядка
$\mu/10$, в сотни раз больше реальных), а движением \textbf{дислокаций} --- линейных дефектов, вдоль которых
одна часть кристалла уже сдвинута относительно другой.

Проще всего представить \textbf{краевую дислокацию}: в кристалл сверху вставлена лишняя полуплоскость атомов,
обрывающаяся на плоскости скольжения (рис.~\ref{fig:f4_dislocation}а). Сдвиг, который несёт дислокация, задаёт
\textbf{вектор Бюргерса} $\mathbf b$; в цирконии для лёгкого скольжения $\mathbf b = \tfrac13\langle11\bar20\rangle$,
$|\mathbf b| = a$. Когда дислокация проходит через кристалл, верхняя половина сдвигается на $\mathbf b$.

\fig{f4_dislocation}{а) Краевая дислокация: над плоскостью скольжения вставлена лишняя полуплоскость (красные
атомы), соседние столбцы раздвинуты. б) Напряжение $\sigma_{xx}$ вокруг неё (изотропная теория упругости): над
плоскостью скольжения --- сжатие, под ней --- растяжение; убывает как $1/r$.}

Вокруг дислокации решётка искажена, и её поле напряжений спадает как $1/r$:
\begin{equation}
\sigma_{xx} = -\frac{\mu b}{2\pi(1-\nu)}\,\frac{y\,(3x^2+y^2)}{(x^2+y^2)^2}.
\end{equation}
Энергия дислокации на единицу длины $\approx \mu b^2/2$ --- несколько эВ на каждое межатомное расстояние. Поэтому
дислокации никогда не находятся в равновесии: они появляются при деформации и остаются, пока их не уберёт отжиг.

\begin{zrbox}
Дислокации важны для гидридов трижды:
\begin{itemize}
\item в растянутой области под краевой дислокацией водороду «просторнее» (глава~\ref{ch:diffusion}), и туда он стекается ---
  дислокации служат местами зарождения гидридов;
\item гидрид при выпадении сам рождает дислокации вокруг себя: его несоответствие слишком велико, чтобы матрица
  выдержала его упруго (глава~\ref{ch:plasticity});
\item эти дислокации остаются после растворения гидрида и «помнят» его место --- при следующем охлаждении гидрид
  выпадает там же. Это \textbf{эффект памяти}.
\end{itemize}
\end{zrbox}

\section{Системы скольжения}

Дислокация скользит в определённой плоскости и в определённом направлении --- вместе они образуют
\textbf{систему скольжения}. Почти всегда это плотнейшая плоскость и плотнейшее направление в ней: вектор Бюргерса
тогда наименьший (энергия $\propto b^2$), а плотные плоскости разнесены дальше всего и сдвигаются легче. В ГЦК
это $\{111\}\langle110\rangle$ --- 12 систем, поэтому ГЦК металлы пластичны в любом направлении. В ГПУ картина беднее.

\begin{center}\small
\begin{tabularx}{\linewidth}{>{\raggedright\arraybackslash}p{3.3cm}>{\raggedright\arraybackslash}p{3.0cm}c>{\raggedright\arraybackslash}X}
\toprule
система & плоскость, направление & систем & в цирконии \\
\midrule
призматическая $\langle a\rangle$ & $\{10\bar10\}\langle11\bar20\rangle$ & 3 & основная, самая лёгкая \\
базисная $\langle a\rangle$ & $(0001)\langle11\bar20\rangle$ & 3 & заметно твёрже призматической, включается при высоких $T$ \\
пирамидальная $\langle a\rangle$ & $\{10\bar11\}\langle11\bar20\rangle$ & 6 & вспомогательная \\
пирамидальная $\langle c+a\rangle$ & $\{10\bar11\}\langle11\bar23\rangle$ & 12 & единственная дающая деформацию вдоль $c$; в разы твёрже призматической \\
двойникование растяжения & $\{10\bar12\}\langle\bar1011\rangle$ & 6 & при растяжении вдоль $c$ (сжатии поперёк) \\
двойникование сжатия & $\{11\bar22\}\langle11\bar2\bar3\rangle$ & 6 & при сжатии вдоль $c$, низкие $T$ \\
\bottomrule
\end{tabularx}
\end{center}

Все системы $\langle a\rangle$ сдвигают атомы только в базисной плоскости: удлинить или укоротить зерно вдоль
оси $c$ они не могут. Произвольное изменение формы зерна требует пяти \emph{независимых} систем (критерий фон
Мизеса), а призматическое скольжение даёт всего две. Недостающее добавляют пирамидальное $\langle c+a\rangle$
скольжение и двойникование --- они «дорогие», поэтому зёрна, нагруженные вдоль $c$, гораздо прочнее нагруженных
поперёк. Отсюда сильная анизотропия циркониевых труб (главы~\ref{ch:texture} и \ref{ch:plasticity}).

\section{Частичные дислокации и дефекты упаковки}

Полная дислокация $\tfrac13\langle11\bar20\rangle$ в плотной плоскости может расщепиться на две
\textbf{частичные дислокации Шокли} типа $\tfrac13\langle10\bar10\rangle$ длиной $a/\sqrt3$: атомы верхнего слоя
переходят из ямки B в ямку B не напрямую через «горб», а через соседнюю ямку C. Энергия $b^2$ при расщеплении
уменьшается ($a^2 \to 2\cdot a^2/3$), но между частичными остаётся полоса, где укладка нарушена, ---
\textbf{дефект упаковки} (рис.~\ref{fig:p4_shockley}б). Его энергия на единицу площади $\gamma_{\text{ДУ}}$
определяет, насколько далеко разойдутся частичные: чем меньше $\gamma_{\text{ДУ}}$, тем шире расщепление.

\fig{p4_shockley}{Плотные слои в ГПУ (вид вдоль $\langle11\bar20\rangle$, позиции A, B, C разнесены по горизонтали).
а) Укладка ABAB. б) Частичная дислокация Шокли прошла между вторым и третьим слоями: все слои выше сдвинуты на
треть периода, укладка ABCACA --- дефект упаковки, внутри которого местная ГЦК последовательность ABC. в) Частичные
на каждом втором слое превращают ABAB в ABCABC --- ГЦК.}

\begin{zrbox}
Именно так перестраивается подрешётка циркония при образовании $\delta$-гидрида. Частичные дислокации Шокли
проходят по каждой второй базисной плоскости, превращая укладку ABAB в ABCABC; одновременно водород заполняет
тетраэдрические пустоты, и решётка разбухает на 17\,\% по объёму. Сдвиги частичных в разных слоях могут идти в
трёх эквивалентных направлениях; если они чередуются, суммарный сдвиг формы почти гасится, и остаётся в
основном расширение --- собственная деформация гидрида (глава~\ref{ch:elastic}). Поэтому ориентационное
соотношение $(0001)_\alpha\parallel\{111\}_\delta$ не случайность, а прямое следствие механизма.
\end{zrbox}

\section{Сколько дефектов в металле}

Плотность дислокаций $\rho$ --- суммарная длина их линий в единице объёма (м/м$^3$ = м$^{-2}$). В хорошо отожжённом
металле $\rho\sim10^{10}$--$10^{12}$~м$^{-2}$, после сильной холодной деформации --- $10^{14}$--$10^{15}$~м$^{-2}$.
Среднее расстояние между дислокациями $\rho^{-1/2}$: от 10~мкм до 30~нм. Холоднодеформированная и снятая от
напряжений оболочка содержит дислокаций на порядки больше, чем рекристаллизованная, и это сказывается на всём ---
от прочности до мест зарождения гидридов.

Размер зерна измеряют на шлифе методом секущих (средняя длина отрезка прямой внутри одного зерна, $\bar\ell$) или
по стандартной шкале ASTM: номер $G$ задаёт число зёрен $N = 2^{G-1}$ на квадратный дюйм при увеличении 100.
Пересчёт в средний диаметр: $\bar d\approx 0.25\cdot2^{-(G-1)/2}$~мм, то есть $G = 10$ --- около 11~мкм, $G = 12$ ---
около 5.5~мкм. Зёрна оболочек циркония --- единицы микрометров.

\section{Межфазные границы}

Граница между матрицей и выделением бывает трёх типов (рис.~\ref{fig:f4_interfaces}).
\begin{itemize}
\item \textbf{Когерентная}: атомные плоскости непрерывно переходят через границу; несоответствие решёток
  снимается упругим искажением. Энергия границы мала (0.05--0.2~Дж/м$^2$), но упругая энергия велика и растёт с
  объёмом выделения.
\item \textbf{Полукогерентная}: несоответствие частично снято \textbf{дислокациями несоответствия}, по одной на
  каждые $b/\delta$ плоскостей, где $\delta = (a_\beta - a_\alpha)/a_\alpha$. Энергия границы выше (0.2--0.5).
\item \textbf{Некогерентная}: решётки никак не сопряжены. Энергия границы велика (0.5--1), зато упругое поле
  почти снимается.
\end{itemize}

\fig{f4_interfaces}{Типы межфазных границ: а) когерентная --- решётка выделения подстроена под матрицу;
б) полукогерентная --- периодические дислокации несоответствия ($\perp$); в) некогерентная.}

\begin{idea}
Выделение выбирает компромисс между энергией границы (растёт как площадь, $\propto r^2$) и упругой энергией
(растёт как объём, $\propto r^3$). Маленький зародыш почти всегда когерентен: граница дешёвая, а упругая энергия
при малом объёме ещё невелика. Большой гидрид теряет когерентность, рождая дислокации.
\end{idea}

\section{Границы зёрен}

Металл трубы --- поликристалл из зёрен в несколько микрометров с разной ориентацией решётки. На стыке двух зёрен
--- \textbf{граница зерна}. При малой разориентации $\theta$ это просто стенка дислокаций, и её энергия растёт по
формуле Рида--Шокли $\gamma = \gamma_m(\theta/\theta_m)\bigl(1-\ln(\theta/\theta_m)\bigr)$. При $\theta\gtrsim15^\circ$
дислокации сливаются, и энергия выходит на плато $\gamma_m\approx0.5$--1~Дж/м$^2$ с провалами у особых
разориентаций (рис.~\ref{fig:f4_gb}).

\begin{figure}[htbp]\centering
\includegraphics[width=0.6\linewidth]{f4_gb}
\caption{Энергия границы зерна от угла разориентации: рост по Риду--Шокли для малоугловых границ, плато для
высокоугловых, провалы у специальных разориентаций (совпадающих узлов) --- условно.}
\label{fig:f4_gb}
\end{figure}

\begin{zrbox}
Границы зёрен --- готовые места для зарождения: часть энергии границы «уходит в дело», когда на ней вырастает
выделение (глава~\ref{ch:nucleation}). Гидриды часто выпадают на границах, и межзёренные гидриды могут быть ориентированы иначе, чем
внутризёренные, привязанные к базисной плоскости своего зерна (Qin и др., 2011; Wang и др., 2019). Это ещё один
путь к радиальным гидридам, когда внутри зёрен им мешает текстура.
\end{zrbox}

\begin{exercise}
Объёмное несоответствие гидрида 17\,\%. Сколько это в линейном несоответствии, если считать его изотропным? Через
сколько межплоскостных расстояний пришлось бы ставить дислокацию несоответствия, чтобы снять его полностью?
\end{exercise}

\begin{exercise}
Оцените энергию дислокации на единицу длины при $\mu = 33$~ГПа, $b = 0.323$~нм. Сколько это в эВ на межатомное
расстояние?
\end{exercise}

\begin{exercise}
Покажите, что расщепление $\tfrac13[11\bar20]\to\tfrac13[10\bar10]+\tfrac13[01\bar10]$ выгодно по критерию $b^2$ (правило
Франка). Найдите длину вектора частичной дислокации для циркония.
\end{exercise}

\begin{exercise}
Холоднодеформированный металл имеет $\rho = 10^{15}$~м$^{-2}$. Какую энергию (Дж/м$^3$ и Дж/моль) запасают дислокации,
если энергия линии $\mu b^2/2$? Сравните с движущей силой выпадения гидрида (несколько кДж на моль H).
\end{exercise}
